\chapterstar{Objetivo General}

Desarrollar un sistema de monitoreo de cultivos de jitomate mediante visión artificial y geolocalización, utilizando un dron equipado con cámara y GPS, una Raspberry Pi~4 y modelos de inteligencia artificial capaces de detectar enfermedades, plagas o anomalías en las plantas, generando notificaciones automáticas al usuario a través de Telegram para facilitar la detección temprana y reducir pérdidas agrícolas.\\

\addcontentsline{toc}{chapter}{Objetivos Específicos}
{\bfseries\fontsize{14pt}{16pt}\selectfont Objetivos Específicos}

\begin{enumerate}[topsep=2pt, itemsep=4pt, parsep=0pt]

\item Automatizar las rutas de vuelo del dron mediante waypoints predefinidos, asegurando la cobertura sistemática y eficiente de las áreas de cultivo designadas para el monitoreo.

\item Establecer la comunicación inalámbrica entre el dron y la Raspberry Pi~4 mediante una red WiFi para la transferencia automática de imágenes capturadas durante el recorrido.

\item Desarrollar un sistema de procesamiento digital de imágenes utilizando OpenCV para preparar y optimizar las fotografías obtenidas del cultivo.

\item Implementar un modelo de inteligencia artificial utilizando TensorFlow Lite capaz de identificar anomalías visibles asociadas a enfermedades o plagas en plantas de jitomate.

\item Desarrollar un sistema de notificaciones remotas mediante un bot de Telegram para enviar alertas automáticas, imágenes y ubicación de anomalías detectadas.

\end{enumerate}

\thispagestyle{empty}
